\section{Omitted derivations for Section \ref{sec:model}}\label{app:modelderivations}

\subsection{Independent Cournot equilibrium}

Fix a project type \(j \in \{I,B\}\), and suppress the type index for notational convenience.
Under independent ownership, the incumbent and the entrant solve
\begin{align}
    \max_{q_I \ge 0}\;& (1-q_I-\sigma q_E)q_I, \\
    \max_{q_E \ge 0}\;& (a-q_E-\sigma q_I)q_E.
\end{align}
The first-order conditions are
\begin{align}
    1-2q_I-\sigma q_E &= 0, \\
    a-2q_E-\sigma q_I &= 0.
\end{align}
Solving the linear system yields
\begin{align}
    q_I^D &= \frac{2-a\sigma}{4-\sigma^2}, \\
    q_E^D &= \frac{2a-\sigma}{4-\sigma^2}.
\end{align}
Substituting these quantities into inverse demand gives
\begin{align}
    p_I^D &= 1-q_I^D-\sigma q_E^D = q_I^D, \\
    p_E^D &= a-q_E^D-\sigma q_I^D = q_E^D.
\end{align}
Hence independent profits are
\begin{align}
    \pi_I^D &= (q_I^D)^2
    = \left(\frac{2-a\sigma}{4-\sigma^2}\right)^2, \\
    \pi_E^D &= (q_E^D)^2
    = \left(\frac{2a-\sigma}{4-\sigma^2}\right)^2.
\end{align}

\subsection{Merged-firm outcome}

Under common ownership, the merged firm chooses \((q_I,q_E)\) to maximize
\begin{equation}
    (1-q_I-\sigma q_E)q_I + (a-q_E-\sigma q_I)q_E + m.
\end{equation}
The first-order conditions are
\begin{align}
    1-2q_I-2\sigma q_E &= 0, \\
    a-2q_E-2\sigma q_I &= 0.
\end{align}
Solving yields
\begin{align}
    q_I^M &= \frac{1-a\sigma}{2(1-\sigma^2)}, \\
    q_E^M &= \frac{a-\sigma}{2(1-\sigma^2)}.
\end{align}
Substituting into joint operating profit yields
\begin{equation}
    \Pi_{\mathrm{op}}^M
    =
    \frac{a^2-2a\sigma+1}{4(1-\sigma^2)}.
\end{equation}
Therefore total merged profit is
\begin{equation}
    \Pi^M
    =
    \frac{a^2-2a\sigma+1}{4(1-\sigma^2)} + m.
\end{equation}
The merger surplus is
\begin{equation}
    \Delta
    =
    \Pi^M-\pi_I^D-\pi_E^D.
\end{equation}

\subsection{Consumer surplus}

A quadratic utility representation that rationalizes the inverse demands in Section \ref{sec:model} is
\begin{equation}
    U(q_I,q_E)
    =
    y + q_I + a q_E
    - \frac{1}{2}
      \left(
        q_I^2 + 2\sigma q_I q_E + q_E^2
      \right),
\end{equation}
where \(y\) is numeraire consumption.
Inverse demands are the marginal utilities:
\begin{align}
    p_I &= \frac{\partial U}{\partial q_I}
    = 1-q_I-\sigma q_E, \\
    p_E &= \frac{\partial U}{\partial q_E}
    = a-q_E-\sigma q_I.
\end{align}
Hence consumer surplus is
\begin{equation}
    C(q_I,q_E)
    =
    \frac{1}{2}
    \left(
        q_I^2 + q_E^2 + 2\sigma q_I q_E
    \right).
\end{equation}
Evaluating this expression at \((q_I^D,q_E^D)\) and \((q_I^M,q_E^M)\) yields \(C^D(j)\) and \(C^M(j)\), respectively.

\subsection{Merger-attempt threshold}

For a successful type-\(j\) project, the joint payoff from attempting a merger is
\begin{equation}
    q\Pi^M(j) + (1-q)\bigl(\pi_I^D(j)+\pi_E^D(j)\bigr) - k.
\end{equation}
The joint payoff from not attempting a merger is
\begin{equation}
    \pi_I^D(j)+\pi_E^D(j).
\end{equation}
Therefore a merger is attempted if and only if
\begin{equation}
    q\Pi^M(j) + (1-q)\bigl(\pi_I^D(j)+\pi_E^D(j)\bigr) - k
    \ge
    \pi_I^D(j)+\pi_E^D(j),
\end{equation}
which simplifies to
\begin{equation}
    q\Delta_j \ge k.
\end{equation}
This yields the attempt indicator
\begin{equation}
    a_j(q)=\mathbf{1}\{q\Delta_j \ge k\}
\end{equation}
used in Sections \ref{sec:subsidy} and \ref{sec:decentralized}.

\subsection*{A.5 Welfare accounting and the primitive benchmark}

The main text distinguishes between the approval probability \(q\) and the attempt decision
\begin{equation}
    a_j(q)=\mathbf{1}\{q\Delta_j\ge k\}.
\end{equation}
This distinction is essential for consistent welfare accounting.

If \(q\Delta_j<k\), then no merger is attempted in private equilibrium. In that case, successful
entry necessarily results in independent competition, so any successful-project welfare object must
coincide with the corresponding independent-ownership value.

If \(q\Delta_j\ge k\), then a merger is attempted. Conditional on attempt, approved acquisition
occurs with probability \(q\), rejected merger with probability \(1-q\), and the real resource cost \(k\)
is incurred regardless of approval. This is why planner-specific successful-project values in the main
text take the piecewise form
\begin{equation}
    B_j^p(q)=B_j^{p,D}+a_j(q)\Bigl[q\bigl(B_j^{p,A}-B_j^{p,D}\bigr)-\kappa^p k\Bigr],
\end{equation}
where \(p\) indexes the planner.

For the primitive benchmark introduced in revised Section \ref{sec:subsidy}, define
\begin{equation}
    V_j^D \equiv C^D(j)+\pi_I^D(j)+\pi_E^D(j),
    \qquad
    V_j^A \equiv C^M(j)+\Pi^M(j),
\end{equation}
and
\begin{equation}
    \Omega_j \equiv V_j^A-V_j^D.
\end{equation}
Using the definition of private merger surplus,
\begin{equation}
    \Delta_j=\Pi^M(j)-\pi_I^D(j)-\pi_E^D(j),
\end{equation}
we can write
\begin{equation}
    \Omega_j
    =
    \bigl(C^M(j)-C^D(j)\bigr)+\Delta_j.
\end{equation}
Hence the primitive benchmark successful-project value is
\begin{equation}
    B_j^{PM}(q)\equiv V_j^D+a_j(q)\bigl[q\Omega_j-k\bigr].
\end{equation}
This benchmark uses only product-market primitives together with the same endogenous merger-attempt
condition that governs private continuation values.
